% Options for packages loaded elsewhere
% Options for packages loaded elsewhere
\PassOptionsToPackage{unicode}{hyperref}
\PassOptionsToPackage{hyphens}{url}
\PassOptionsToPackage{dvipsnames,svgnames,x11names}{xcolor}
%
\documentclass[
  letterpaper,
  DIV=11,
  numbers=noendperiod]{scrartcl}
\usepackage{xcolor}
\usepackage{amsmath,amssymb}
\setcounter{secnumdepth}{-\maxdimen} % remove section numbering
\usepackage{iftex}
\ifPDFTeX
  \usepackage[T1]{fontenc}
  \usepackage[utf8]{inputenc}
  \usepackage{textcomp} % provide euro and other symbols
\else % if luatex or xetex
  \usepackage{unicode-math} % this also loads fontspec
  \defaultfontfeatures{Scale=MatchLowercase}
  \defaultfontfeatures[\rmfamily]{Ligatures=TeX,Scale=1}
\fi
\usepackage{lmodern}
\ifPDFTeX\else
  % xetex/luatex font selection
\fi
% Use upquote if available, for straight quotes in verbatim environments
\IfFileExists{upquote.sty}{\usepackage{upquote}}{}
\IfFileExists{microtype.sty}{% use microtype if available
  \usepackage[]{microtype}
  \UseMicrotypeSet[protrusion]{basicmath} % disable protrusion for tt fonts
}{}
\makeatletter
\@ifundefined{KOMAClassName}{% if non-KOMA class
  \IfFileExists{parskip.sty}{%
    \usepackage{parskip}
  }{% else
    \setlength{\parindent}{0pt}
    \setlength{\parskip}{6pt plus 2pt minus 1pt}}
}{% if KOMA class
  \KOMAoptions{parskip=half}}
\makeatother
% Make \paragraph and \subparagraph free-standing
\makeatletter
\ifx\paragraph\undefined\else
  \let\oldparagraph\paragraph
  \renewcommand{\paragraph}{
    \@ifstar
      \xxxParagraphStar
      \xxxParagraphNoStar
  }
  \newcommand{\xxxParagraphStar}[1]{\oldparagraph*{#1}\mbox{}}
  \newcommand{\xxxParagraphNoStar}[1]{\oldparagraph{#1}\mbox{}}
\fi
\ifx\subparagraph\undefined\else
  \let\oldsubparagraph\subparagraph
  \renewcommand{\subparagraph}{
    \@ifstar
      \xxxSubParagraphStar
      \xxxSubParagraphNoStar
  }
  \newcommand{\xxxSubParagraphStar}[1]{\oldsubparagraph*{#1}\mbox{}}
  \newcommand{\xxxSubParagraphNoStar}[1]{\oldsubparagraph{#1}\mbox{}}
\fi
\makeatother


\usepackage{longtable,booktabs,array}
\usepackage{calc} % for calculating minipage widths
% Correct order of tables after \paragraph or \subparagraph
\usepackage{etoolbox}
\makeatletter
\patchcmd\longtable{\par}{\if@noskipsec\mbox{}\fi\par}{}{}
\makeatother
% Allow footnotes in longtable head/foot
\IfFileExists{footnotehyper.sty}{\usepackage{footnotehyper}}{\usepackage{footnote}}
\makesavenoteenv{longtable}
\usepackage{graphicx}
\makeatletter
\newsavebox\pandoc@box
\newcommand*\pandocbounded[1]{% scales image to fit in text height/width
  \sbox\pandoc@box{#1}%
  \Gscale@div\@tempa{\textheight}{\dimexpr\ht\pandoc@box+\dp\pandoc@box\relax}%
  \Gscale@div\@tempb{\linewidth}{\wd\pandoc@box}%
  \ifdim\@tempb\p@<\@tempa\p@\let\@tempa\@tempb\fi% select the smaller of both
  \ifdim\@tempa\p@<\p@\scalebox{\@tempa}{\usebox\pandoc@box}%
  \else\usebox{\pandoc@box}%
  \fi%
}
% Set default figure placement to htbp
\def\fps@figure{htbp}
\makeatother





\setlength{\emergencystretch}{3em} % prevent overfull lines

\providecommand{\tightlist}{%
  \setlength{\itemsep}{0pt}\setlength{\parskip}{0pt}}



 


\KOMAoption{captions}{tableheading}
\makeatletter
\@ifpackageloaded{caption}{}{\usepackage{caption}}
\AtBeginDocument{%
\ifdefined\contentsname
  \renewcommand*\contentsname{Table of contents}
\else
  \newcommand\contentsname{Table of contents}
\fi
\ifdefined\listfigurename
  \renewcommand*\listfigurename{List of Figures}
\else
  \newcommand\listfigurename{List of Figures}
\fi
\ifdefined\listtablename
  \renewcommand*\listtablename{List of Tables}
\else
  \newcommand\listtablename{List of Tables}
\fi
\ifdefined\figurename
  \renewcommand*\figurename{Figure}
\else
  \newcommand\figurename{Figure}
\fi
\ifdefined\tablename
  \renewcommand*\tablename{Table}
\else
  \newcommand\tablename{Table}
\fi
}
\@ifpackageloaded{float}{}{\usepackage{float}}
\floatstyle{ruled}
\@ifundefined{c@chapter}{\newfloat{codelisting}{h}{lop}}{\newfloat{codelisting}{h}{lop}[chapter]}
\floatname{codelisting}{Listing}
\newcommand*\listoflistings{\listof{codelisting}{List of Listings}}
\makeatother
\makeatletter
\makeatother
\makeatletter
\@ifpackageloaded{caption}{}{\usepackage{caption}}
\@ifpackageloaded{subcaption}{}{\usepackage{subcaption}}
\makeatother
\usepackage{bookmark}
\IfFileExists{xurl.sty}{\usepackage{xurl}}{} % add URL line breaks if available
\urlstyle{same}
\hypersetup{
  pdftitle={Redes Neuronales Convolucionales},
  pdfauthor={Mariana Rodríguez},
  colorlinks=true,
  linkcolor={blue},
  filecolor={Maroon},
  citecolor={Blue},
  urlcolor={Blue},
  pdfcreator={LaTeX via pandoc}}


\title{Redes Neuronales Convolucionales}
\author{Mariana Rodríguez}
\date{}
\begin{document}
\maketitle


\section{Introducción}\label{introducciuxf3n}

Una \textbf{red neuronal convolucional} (CNN, \emph{Convolutional Neural
Network}) es una arquitectura de aprendizaje profundo diseñada para
procesar datos con estructura de tipo rejilla, como imágenes.

Formalmente, una CNN se define como una composición de funciones:

\[
f(x; \theta) = f^{(L)} \circ f^{(L-1)} \circ \cdots \circ f^{(1)}(x),
\]

donde \(x \in \mathbb{R}^{H \times W \times C}\) es la entrada y
\(\theta\) representa los parámetros del modelo.

\section{Operación de Convolución}\label{operaciuxf3n-de-convoluciuxf3n}

Sea una imagen:

\[
X \in \mathbb{R}^{H \times W}
\]

y un kernel:

\[
K \in \mathbb{R}^{k_h \times k_w}.
\]

La salida de la convolución discreta está dada por:

\[
Y(i,j) = \sum_{u=0}^{k_h-1} \sum_{v=0}^{k_w-1} K(u,v)\, X(i+u, j+v).
\]

En aprendizaje profundo, esta operación corresponde técnicamente a una
correlación cruzada.

\section{Convolución Multicanal}\label{convoluciuxf3n-multicanal}

Para una entrada con \(C\) canales:

\[
X \in \mathbb{R}^{H \times W \times C}, \quad
K \in \mathbb{R}^{k_h \times k_w \times C},
\]

la convolución se define como:

\[
Y(i,j) = \sum_{c=1}^{C} \sum_{u=0}^{k_h-1} \sum_{v=0}^{k_w-1} K(u,v,c)\, X(i+u, j+v, c).
\]

\section{Múltiples Filtros}\label{muxfaltiples-filtros}

Si se utilizan \(F\) filtros:

\[
K^{(f)} \in \mathbb{R}^{k_h \times k_w \times C}, \quad f = 1,\dots,F,
\]

entonces la salida es:

\[
Y \in \mathbb{R}^{H' \times W' \times F},
\]

donde:

\[
Y_f(i,j) = \sum_{c,u,v} K^{(f)}(u,v,c)\, X(i+u, j+v, c).
\]

\section{Stride y Padding}\label{stride-y-padding}

\subsection{Stride}\label{stride}

El stride \(s\) controla el desplazamiento del kernel. Las dimensiones
de salida son:

\[
H' = \left\lfloor \frac{H - k_h + 2p}{s} \right\rfloor + 1, \quad
W' = \left\lfloor \frac{W - k_w + 2p}{s} \right\rfloor + 1.
\]

\subsection{Padding}\label{padding}

Se define como la adición de ceros alrededor de la imagen:

\[
X_{\text{pad}} \in \mathbb{R}^{(H+2p)\times(W+2p)\times C}.
\]

\section{Función de Activación}\label{funciuxf3n-de-activaciuxf3n}

Después de la convolución se aplica una función no lineal. Un ejemplo
común es la función ReLU:

\[
\phi(x) = \max(0,x).
\]

La salida es:

\[
Z = \phi(Y).
\]

\section{Pooling}\label{pooling}

El pooling reduce la dimensionalidad.

\subsection{Max Pooling}\label{max-pooling}

\[
Y(i,j) = \max_{(u,v)\in \mathcal{R}} X(i+u, j+v).
\]

\subsection{Average Pooling}\label{average-pooling}

\[
Y(i,j) = \frac{1}{|\mathcal{R}|} \sum_{(u,v)\in \mathcal{R}} X(i+u, j+v).
\]

\section{Capa Totalmente Conectada}\label{capa-totalmente-conectada}

Se transforma el tensor en un vector:

\[
X \in \mathbb{R}^{H \times W \times C} \rightarrow x \in \mathbb{R}^{n}.
\]

Luego:

\[
y = Wx + b,
\]

donde \(W \in \mathbb{R}^{m \times n}\) y \(b \in \mathbb{R}^{m}\).

\section{Función de Pérdida}\label{funciuxf3n-de-puxe9rdida}

\subsection{Softmax}\label{softmax}

\[
\hat{y}_i = \frac{e^{z_i}}{\sum_{j} e^{z_j}}.
\]

\subsection{Entropía Cruzada}\label{entropuxeda-cruzada}

\[
\mathcal{L} = -\sum_{i} y_i \log(\hat{y}_i).
\]

\section{Entrenamiento}\label{entrenamiento}

El problema de optimización es:

\[
\min_{\theta} \mathcal{L}(f(x;\theta), y).
\]

Se utiliza descenso de gradiente:

\[
\theta \leftarrow \theta - \eta \nabla_{\theta} \mathcal{L}.
\]

\section{Propiedades Fundamentales}\label{propiedades-fundamentales}

\subsection{Conectividad Local}\label{conectividad-local}

Cada neurona depende de una región pequeña de la entrada.

\subsection{Compartición de
Parámetros}\label{comparticiuxf3n-de-paruxe1metros}

El mismo kernel se aplica en toda la imagen.

\subsection{Invarianza Traslacional}\label{invarianza-traslacional}

El modelo detecta patrones independientemente de su posición.

\section{Flujo de una CNN}\label{flujo-de-una-cnn}

\[
\text{Imagen} \rightarrow \text{Convolución} \rightarrow \text{ReLU} \rightarrow \text{Pooling} \rightarrow \cdots \rightarrow \text{Fully Connected} \rightarrow \text{Softmax}
\]

\section{Interpretación
Matemática}\label{interpretaciuxf3n-matemuxe1tica}

Una CNN puede verse como una función:

\[
f: \mathbb{R}^{H \times W \times C} \to \mathbb{R}^{k},
\]

compuesta por operadores lineales, no lineales y de reducción.

\section{Conclusión}\label{conclusiuxf3n}

Las redes neuronales convolucionales son modelos altamente eficientes
para el procesamiento de imágenes debido a la combinación de
convoluciones, no linealidades y reducción de dimensionalidad, junto con
propiedades estructurales como la compartición de parámetros y la
conectividad local.




\end{document}
